\documentclass[12pt,a4paper]{article}
\usepackage[utf8]{inputenc}
\usepackage[vietnamese]{babel}
\usepackage{amsmath,amssymb,amsthm}
\usepackage{geometry}
\usepackage{enumitem}
\usepackage[most]{tcolorbox}
\usepackage{hyperref}
\usepackage{fancyhdr}
\usepackage{tikz}
\usepackage{eso-pic}
\usepackage{xcolor}
\geometry{a4paper, margin=2cm, bottom=2.5cm}
\hypersetup{colorlinks=true, linkcolor=blue!70!black}
\setlist[itemize]{topsep=4pt, itemsep=2pt}
\setlist[enumerate]{topsep=4pt, itemsep=2pt}
\AddToShipoutPictureFG{%
  \AtPageCenter{%
    \begin{tikzpicture}[remember picture, overlay]
      \node[rotate=45, scale=1, opacity=0.06]{\fontsize{60}{72}\selectfont\bfseries\sffamily Đặng Tấn Quí};
    \end{tikzpicture}%
  }%
}
\pagestyle{fancy}
\fancyhf{}
\fancyhead[L]{\small\textit{QUY HOẠCH TUYẾN TÍNH}}
\fancyhead[R]{\small\textit{Đặng Tấn Quí}}
\fancyfoot[C]{\thepage}
\fancyfoot[L]{\footnotesize\textcopyright\ Đặng Tấn Quí}
\fancyfoot[R]{\footnotesize SĐT: 0377\,122\,210}
\renewcommand{\headrulewidth}{0.4pt}
\renewcommand{\footrulewidth}{0.4pt}
\fancypagestyle{plain}{\fancyhf{}\fancyfoot[C]{\thepage}\fancyfoot[L]{\footnotesize\textcopyright\ Đặng Tấn Quí}\fancyfoot[R]{\footnotesize SĐT: 0377\,122\,210}\renewcommand{\headrulewidth}{0pt}\renewcommand{\footrulewidth}{0.4pt}}
\newtcolorbox{dinhnghia}[1][]{enhanced, breakable, colback=blue!3!white, colframe=blue!60!black, fonttitle=\bfseries, title={Định nghĩa}, #1}
\newtcolorbox{dinhly}[1][]{enhanced, breakable, colback=green!3!white, colframe=green!50!black, fonttitle=\bfseries, title={Định lý}, #1}
\newtcolorbox{tinhchat}[1][]{enhanced, breakable, colback=yellow!5!white, colframe=orange!50!black, fonttitle=\bfseries, title={Tính chất}, #1}
\newtcolorbox{phuongphap}[1][]{enhanced, breakable, colback=gray!5!white, colframe=gray!50!black, fonttitle=\bfseries, title={Phương pháp}, #1}
\newtcolorbox{luuy}[1][]{enhanced, breakable, colback=red!3!white, colframe=red!50!black, fonttitle=\bfseries, title={Lưu ý}, #1}
\newtcolorbox{congthuc}[1][]{enhanced, breakable, colback=cyan!3!white, colframe=cyan!50!black, fonttitle=\bfseries, title={Công thức}, #1}
\newtcolorbox{vidu}[1][]{enhanced, breakable, colback=violet!3!white, colframe=violet!50!black, fonttitle=\bfseries, title={Ví dụ}, #1}

\begin{document}
\thispagestyle{empty}
\begin{tikzpicture}[remember picture, overlay]
  \fill[blue!80!black] (current page.south west) rectangle (current page.north east);
  \node[anchor=west, white, font=\fontsize{26}{32}\bfseries\selectfont]
    at ([xshift=3cm, yshift=-7.5cm]current page.north west) {QUY HOẠCH TUYẾN TÍNH};
  \node[anchor=west, white, font=\fontsize{18}{24}\selectfont]
    at ([xshift=3cm, yshift=-9cm]current page.north west) {Simplex, đối ngẫu, vận tải, phân công};
  \node[anchor=west, white, font=\Large\bfseries] at ([xshift=3cm, yshift=-13cm]current page.north west) {Biên soạn:};
  \node[anchor=west, yellow!80!white, font=\fontsize{22}{28}\bfseries\selectfont] at ([xshift=3cm, yshift=-14.5cm]current page.north west) {ĐẶNG TẤN QUÍ};
  \node[anchor=south east, white, font=\Large\bfseries] at ([xshift=-2cm, yshift=3cm]current page.south east) {\the\year};
\end{tikzpicture}
\newpage
\tableofcontents
\newpage
%% ================================================================
\section{Bài toán quy hoạch tuyến tính}

\begin{dinhnghia}[title={Dạng tổng quát}]
  $\min\; \mathbf{c}^T\mathbf{x}$ \quad t.c. $\mathbf{A}\mathbf{x} \le \mathbf{b}$, $\mathbf{x} \ge \mathbf{0}$.
\end{dinhnghia}

\begin{dinhnghia}[title={Dạng chuẩn và chính tắc}]
  \textbf{Dạng chính tắc}: $\min\;\mathbf{c}^T\mathbf{x}$ t.c. $\mathbf{A}\mathbf{x} = \mathbf{b}$, $\mathbf{x} \ge \mathbf{0}$.

  Chuyển đổi: thêm biến bù $s \ge 0$: $\mathbf{A}\mathbf{x} + \mathbf{s} = \mathbf{b}$.
\end{dinhnghia}

\begin{dinhnghia}[title={Nghiệm cơ bản khả thi (BFS)}]
  Chọn $m$ cột độc lập $\mathbf{B}$ (cơ sở), $\mathbf{x}_B = \mathbf{B}^{-1}\mathbf{b} \ge \mathbf{0}$, $\mathbf{x}_N = \mathbf{0}$.
\end{dinhnghia}

\begin{dinhly}
  Nếu bài toán QHTT có nghiệm tối ưu hữu hạn thì tồn tại BFS tối ưu (đỉnh của đa diện).
\end{dinhly}

%% ================================================================
\section{Phương pháp đơn hình (Simplex)}

\begin{phuongphap}[title={Thuật toán Simplex}]
  \begin{enumerate}
    \item Xuất phát từ BFS $\mathbf{x}_B = \mathbf{B}^{-1}\mathbf{b}$.
    \item Tính chi phí rút gọn: $\bar{c}_j = c_j - \mathbf{c}_B^T \mathbf{B}^{-1}\mathbf{a}_j$.
    \item Nếu $\bar{c}_j \ge 0$ $\forall j$: tối ưu. Ngừng.
    \item Chọn $j$ vào cơ sở: $\bar{c}_j < 0$ (quy tắc Bland, Dantzig, \ldots).
    \item Tỉ số tối thiểu: chọn $i$ ra khỏi cơ sở (pivot).
    \item Cập nhật bảng, lặp.
  \end{enumerate}
\end{phuongphap}

\begin{phuongphap}[title={Pha 1 --- Tìm BFS ban đầu}]
  Thêm biến giả $\mathbf{y}$: $\min \sum y_i$ t.c. $\mathbf{A}\mathbf{x} + \mathbf{y} = \mathbf{b}$, $\mathbf{x}, \mathbf{y} \ge 0$. Nếu $\min = 0$: tìm được BFS cho bài toán gốc.
\end{phuongphap}

\begin{luuy}
  Simplex: số bước tồi nhất là mũ, nhưng trung bình rất hiệu quả. Phương pháp điểm trong (interior-point) có độ phức tạp đa thức.
\end{luuy}

%% ================================================================
\section{Đối ngẫu}

\begin{dinhnghia}[title={Bài toán đối ngẫu}]
  Gốc (P): $\min\;\mathbf{c}^T\mathbf{x}$ t.c. $\mathbf{A}\mathbf{x} \ge \mathbf{b}$, $\mathbf{x} \ge 0$.

  Đối ngẫu (D): $\max\;\mathbf{b}^T\mathbf{y}$ t.c. $\mathbf{A}^T\mathbf{y} \le \mathbf{c}$, $\mathbf{y} \ge 0$.
\end{dinhnghia}

\begin{dinhly}[title={Đối ngẫu yếu}]
  Nếu $\mathbf{x}$ khả thi (P) và $\mathbf{y}$ khả thi (D) thì $\mathbf{c}^T\mathbf{x} \ge \mathbf{b}^T\mathbf{y}$.
\end{dinhly}

\begin{dinhly}[title={Đối ngẫu mạnh}]
  Nếu (P) có nghiệm tối ưu hữu hạn thì (D) cũng vậy và $\mathbf{c}^T\mathbf{x}^* = \mathbf{b}^T\mathbf{y}^*$.
\end{dinhly}

\begin{dinhly}[title={Bù lỏng (Complementary Slackness)}]
  $x_j^*(\mathbf{c}^T - \mathbf{y}^{*T}\mathbf{A})_j = 0$ $\forall j$ và $y_i^*(\mathbf{A}\mathbf{x}^* - \mathbf{b})_i = 0$ $\forall i$.
\end{dinhly}

%% ================================================================
\section{Bài toán vận tải và phân công}

\begin{dinhnghia}[title={Bài toán vận tải}]
  $\min \sum_{i,j} c_{ij}x_{ij}$ t.c. $\sum_j x_{ij} = a_i$, $\sum_i x_{ij} = b_j$, $x_{ij} \ge 0$. Cân bằng: $\sum a_i = \sum b_j$.
\end{dinhnghia}

\begin{phuongphap}[title={Giải bài toán vận tải}]
  \begin{enumerate}
    \item Tìm nghiệm khả thi ban đầu: góc Tây Bắc, chi phí nhỏ nhất, Vogel.
    \item Kiểm tra tối ưu: tính thế vị $u_i, v_j$, kiểm tra $c_{ij} - u_i - v_j \ge 0$.
    \item Cải thiện: tìm ô có $\Delta_{ij} < 0$, lập vòng, chuyển.
  \end{enumerate}
\end{phuongphap}

\begin{dinhnghia}[title={Bài toán phân công}]
  Trường hợp đặc biệt: $a_i = b_j = 1$, $x_{ij} \in \{0, 1\}$. Giải bằng thuật toán Hungary: $O(n^3)$.
\end{dinhnghia}

\end{document}